\documentclass[11pt]{article}
\usepackage[margin=1in]{geometry}
\usepackage{amsmath,amssymb}
\usepackage{booktabs}
\usepackage[hidelinks]{hyperref}
\title{Legacy record: the uncorrected seed-3--9 models and superseded results}
\author{Working report (supplement)}
\date{2026-07-27}

% Provenance record only. Documents a training-time regularisation-coefficient issue in the
% first batch of retrained networks, the resulting model characteristics, and the results that
% were computed on those models and subsequently superseded. The main report
% (methods_results_report.tex) uses the corrected models throughout.

\begin{document}
\maketitle

\begin{abstract}
This supplement documents a set of networks (seeds $3$--$9$, sizes $128$ and $256$) trained
under an unintended regularisation coefficient, and the analysis outputs derived from them,
which have been superseded by re-trained (``corrected'') models. It exists for provenance:
the main report analyses the corrected $256$-unit seeds $3$--$9$ (plus released seeds $0$--$2$)
and uses only the released versions at the other network sizes. No conclusions are drawn here.
\end{abstract}

\section{The regularisation-coefficient issue}
A training script specified the L1 activity ($\ell_{1h}$) and L1 weight ($\ell_{1w}$) penalties
by an exponent $e$ that was mapped to a coefficient by $10^{e}$ with no guard. The value
intended to disable the penalty, $e=0$, produced $10^{0}=1$ rather than $0$; the networks were
therefore trained with an L1 penalty of magnitude $1$ on the mean absolute activity and on the
weights (versus the intended $0$; the L2 coefficients, $10^{-6}$, were unaffected). The other
training scripts in the codebase guard this mapping as $\min(-e,10^{e})$ (so $e=0\mapsto0$,
$e=-6\mapsto10^{-6}$); the affected script did not. It was corrected to use the same guard and
seeds $3$--$9$ were re-trained.

\section{Characteristics of the uncorrected models ($256$-unit)}
\begin{center}
\begin{tabular}{lcc}
\toprule
quantity & uncorrected seeds $3$--$9$ & corrected seeds $3$--$9$ \\
\midrule
$\ell_{1h}$, $\ell_{1w}$ coefficients      & $1$, $1$                & $0$, $0$ \\
L2 coefficients                            & $10^{-6}$               & $10^{-6}$ \\
active units (variance $>10^{-3}$)         & $\approx 71$            & $212$--$220$ \\
operating scale ($20\%$ of active)         & $14$--$15$              & $42$--$44$ \\
mean accuracy (worst task)                 & $0.91$--$0.94$          & $0.94$--$0.96$ \\
mean accuracy (across tasks)               & $0.92$--$0.96$          & $\approx 0.975$ \\
\bottomrule
\end{tabular}
\end{center}
The uncorrected models had roughly one third as many active units, and correspondingly a
smaller operating group number, because the large L1 activity penalty drove most units to zero
activity.

\section{Superseded analysis outputs}
The following outputs were computed on the uncorrected models and are superseded by the
corresponding corrected-model results in the main report. The direction of every effect was
preserved after correction; the magnitudes changed. Representative examples ($256$-unit,
$k=20$ or operating scale as applicable):

\begin{center}
\begin{tabular}{lccl}
\toprule
quantity & uncorrected & corrected & main-report section \\
\midrule
motif-direction containment score (median)     & $+2.52$ & $+7.51$ & containment \\
quotient prediction $A_{\rm cond}$ (median)     & (n/a)\,$^{\dagger}$ & $0.816$ & quotient reduced model \\
mode-concentration PR conn vs.\ random ($10/10$) & preserved & preserved & structural bottleneck \\
rotation-readout ring $>$ matched               & preserved & $104/109$ & rotation \\
\bottomrule
\end{tabular}
\end{center}
{\footnotesize $^{\dagger}$ The quotient-prediction, eigenvalue-matched, taxonomy, and reuse
analyses were recomputed from scratch on the corrected models; per-seed uncorrected values were
not carried forward except where a direct before/after comparison was recorded (e.g.\ the
containment score).}

\section{Storage}
The uncorrected checkpoints and their derived structural outputs (activity clusters, signed-WL
sweeps for $256$-unit seeds $3$--$9$) were moved to a backup directory
(\texttt{results/\_buggy\_l1\_backup\_20260726/}) prior to re-running the pipeline on the
corrected models. The uncorrected fixed-point trees were overwritten in place by the
corrected-model re-run.

\section{Other network sizes}
Networks at $128$, $512$ and $1024$ units were also trained for seeds $3$--$9$ under the same
uncorrected setting; the $128$-unit set was subsequently re-trained. In the main report the
non-$256$ sizes are represented only by the released seeds $0$--$2$; the retrained non-$256$
seeds are not analysed there.

\end{document}
